\documentclass[twoside]{article}
\usepackage{graphicx}
\usepackage{polski}
\usepackage{fancyhdr}
\usepackage{lastpage}
\usepackage{epigraph}
\usepackage{listings}
\usepackage{soul}
\usepackage{color}
\usepackage{multirow}
\usepackage{subcaption}
\usepackage{fancyvrb}
\usepackage{hyperref}
\usepackage{minted}


\title{AiSD Ex. 5 (Pon., Wt.)}
\author{\hl{Mateusz Kamieniecki}}
\date{}

\lstset{
    numbers=left
}

\makeatletter

\begin{document}

\pagestyle{fancy}
\fancyhead{}
\fancyhead[L]{\@title}
\fancyhead[R]{\@author}
\fancyhead[C]{\@date}

\cfoot{\thepage\ / \pageref{LastPage}}


\newpage
\section{Zadania}


\subsection{Zadanie 0}
Proszę (już tradycyjnie) uporządkować swoje repozytorium, aby organizacja projektów wyglądała jak poniżej. 
Oczywiście w tej strukturze mogą znajdować się również inne katalogi i pliki, ale proszę
dokładnie odwzorować przedstawiony porządek projektów.

\label{struktura}
\begin{verbatim}
.korzeń_repozytorium
+-- CW_0
|   +-- wszystko, co było
|
+-- CW_1
|   +-- AiSD2026ZEx1
|       +-- pom.xml
|
+-- CW_2
|   +-- AiSD2026ZEx2
|       +-- pom.xml
|
+-- CW_3
|   +-- AiSD2026ZEx3
|       +-- pom.xml
|
+-- CW_4
|   +-- AiSD2026ZEx4
|       +-- pom.xml
|
+-- CW_5
    +-- AiSD2026ZEx5
    |   +-- pom.xml
    |
    +-- README.md
\end{verbatim}

\newpage

\subsection{Zadanie 1}
Proszę zaimplementować program pozwalający na wykonywanie bezstratnych kompresji i dekompresji pojedynczych plików.

\subsubsection{Struktura}

Chcąc osiągnąć ten cel, należy stworzyć projekt \textit{Maven-owy} o parametrach strukturalnych:
\begin{itemize}
    \item \textbf{groupId:} \verb|pl.edu.pw.ee|;
    \item \textbf{artifactId:} \verb|AiSD2026ZEx5|.
\end{itemize}

Powyższe informacje powinny Państwu sugerować, że struktura tego projektu będzie analogiczna 
do wszystkiego, co do tej pory robiliśmy.

\subsubsection{Kompilacja + uruchomienie}

Projekt \textit{Maven-owy} powinien pozwolić na wykonanie fazy \textbf{package} 
(najlepiej w~wariancie z wiersza poleceń, komendą: \verb|mvn clean package|), której wynikiem
powinno być stworzenie uruchamialnego archiwum JAR: \verb|AiSD2026ZEx5.jar|.\\

Przykładowe uruchomienie mogłoby wyglądać następująco:
\begin{minted}[breaklines]{bash}
java -jar AiSD2026ZEx5.jar -m comp -s sciezka/do/katalogu/niemanie.txt -d sciezka/do/innego/katalogu/niemanie_txt.comp -l 2
\end{minted}

w którym argument:
\begin{itemize}
    \item \verb|-m| może przyjmować jedną z dwóch wartości: \verb|comp|, \verb|decomp| -- w zależności od trybu działania programu, odpowiednio dla kompresji i dekompresji;
    \item \verb|-s| określa ścieżkę do pliku wejściowego (poddawanego kompresji lub dekompresji);
    \item \verb|-d| określa ścieżkę do pliku wyjściowego, w którym powinien zostać zapisany wynik działania programu (wynikowy plik po kompresji lub dekompresji);
    \item \verb|-l| określa maksymalną długość łańcuchów, które powinny stanowić pojedynczą sekwencję odpowiadającą słowu kodowemu (np. \verb|-l 2| oznacza, że zamiast wyznaczać słowa dla pojedynczych znaków, program będzie kodował pary znaków). Domyślna wartość tego argumentu wynosi \verb|l=1|.
\end{itemize}

\subsection{Pozostałe wymagania}
\begin{enumerate}
    \item Państwa implementacje powinny być przetestowane zarówno testami jednostkowymi (polecam\
    \verb|JaCoCo| jako wyszukiwanie miejsc, które wymagają Państwa uwagi) jak i testami akceptacyjnymi (czyli użytkowo).
    \item Program powinien być odporny na błędy (np.: podanie nieprawidłowych argumentów, podanie ścieżki do nieistniejącego pliku wejściowego).
    \item Program nie powinien wymagać dodatkowej interakcji z użytkownikiem w~trakcie działania.
    \item Sytuacje wyjątkowe (zgłoszenie wyjątków i błędów) powinno jawnie pochodzić z Państwa kodu.\
    Inaczej: jeśli program zgłosi jakiś wyjątek, to ten wyjątek ma być przez Państwa jawnie (literalnie) zapisany w kodzie.
    \item Wynikowy program (wraz z przeprowadzeniem kompilacji) musi poprawnie działać na komputerach laboratoryjnych.
    \item Można korzystać ze struktur danych i abstrakcyjnych typów danych, które oferuje interfejs \verb|List|.\
    Można również korzystać z kodu, który pojawił się na wszystkich poprzednich ćwiczeniach.\
    Jeśli potrzebują Państwo czegoś więcej, to trzeba to samodzielnie napisać.\
    \textbf{W skrócie:} nie powinni Państwo wykorzystywać (poza pakietem testów) dodatkowych bibliotek.
\end{enumerate}

\subsubsection{A jeśli...}
\begin{enumerate}
    \item Nie będę weryfikował istnienia testów.\
    Jednak zawsze warto mieć testy (to taka dobra praktyka inżynierska).
    \item Można rozszerzyć działanie programu. Można mieć dodatkowe parametry uruchomienia.\
    Ale... podane w poleceniu argumenty (w tej formie, z taką nazwą) mają być przez Państwa program obsłużone prawidłowo.
\end{enumerate}

\subsection{README}
W pliku \verb|README.md| (umieszczonym w miejscu zgodnym ze strukturą (patrz: punkt \ref{struktura})
należy zawrzeć instrukcję pozwalającą na kompilację i uruchomienie kodu.
Zależy mi na działającej, przykładowej instrukcji, którą mogę wziąć i bez wprowadzania zmian uruchomić.

Jeśli będą Państwo mieli do przekazania jakiś komentarz odnośnie swojej pracy, to plik 
\verb|README.md| jest do tego odpowiednim miejscem.


\section{Zakończenie pracy}
Na koniec zajęć (\underline{bez przekroczenia} z góry ustalonego terminu), gdy już będą Państwo
przekonani, że wszystko jest zrobione zgodnie z planem, proszę wysłać swoje zmiany do repozytorium zdalnego, a wersję, która ma być sprawdzona, proszę oznaczyć tagiem (jak zawsze: 
\underline{ze znacznikiem czasowym}) \verb|CW_5_FINAL|.

Tag musi być dla mnie widoczny (w repozytorium zdalnym).
Brak oznaczenia pracy tagiem będzie skutkował niezaliczeniem ćwiczenia 5.

\textbf{Uwaga!} Nie ma obowiązku wykonania opisanego wyżej programu w ramach aktualnego ćwiczenia.
Ale koniecznie muszą Państwo umieścić w repozytorium tag, nawet jeżeli niczego Państwo nie zaimplementowali.

\end{document}
